\documentclass[11pt]{article}
\usepackage[margin=1in]{geometry}
\usepackage{amsmath,amssymb,amsthm}
\usepackage[hidelinks]{hyperref}
\setlength{\emergencystretch}{2em}

\newtheorem*{theorem}{Literature Answer}

\title{Negative Literature Answer to the $V(T)$-Geometry Question\\
in arXiv:1706.08682}
\author{FA/Banach run packet}
\date{21 July 2026}

\begin{document}
\maketitle

\section*{Status}

This is a \texttt{literature\_already\_answered} packet.  It records that the
open question in Hernández--Semenov--Tradacete \cite{HST17} was answered
negatively by the same authors in the later paper \cite{HST20}.  No novelty is
claimed for this note.

\section*{Original question}

For an operator $T:L_\infty\to L_1$, set
\[
 V(T)=\left\{\left(\frac1p,\frac1q\right):
 T:L_p\to L_q\text{ is strictly singular and non-compact}\right\}.
\]
The examples in \cite{HST17} had $V(T)$ assembled from line segments that
were horizontal, vertical, or parallel to the diagonal.  Immediately after
the displayed example, the authors asked whether $V(T)$ must always look
like that.

\section*{Answer}

The later paper \cite{HST20} explicitly restates this question and announces
a negative answer.  Its main realization theorem gives considerably more.

Let $\mathfrak L$ be the family of line segments $\ell\subset(0,1)^2$
such that either
\begin{enumerate}
\item $\ell$ is horizontal or vertical and avoids the central forbidden
region specified in \cite{HST20}; or
\item $\ell$ has positive slope and lies entirely below the diagonal.
\end{enumerate}
For such a segment, let $L_\ell$ be the side of the square selected in
\cite{HST20}.

\begin{theorem}[Hernández--Semenov--Tradacete, Theorem 9]
For every $\ell\in\mathfrak L$, there is an operator
$T:L_\infty\to L_1$ such that
\[
 L(T)=S(T)=L_\ell\cup\ell,\qquad
 K(T)=L_\ell,\qquad
 V(T)=\ell.
\]
Consequently, the question in \cite{HST17} has a negative answer.
\end{theorem}

Indeed, choose an admissible segment below the diagonal of slope $2$ (or
any positive slope other than $1$).  The theorem produces an operator with
$V(T)=\ell$.  This segment is neither horizontal, vertical, nor parallel to
the diagonal, so it cannot have the geometry suggested by the 2017 example.

\section*{Construction and proof idea}

For $0<\alpha,\lambda<1$, let $\Omega_\alpha\subset(0,1)$ be an Ahlfors
$\alpha$-regular set and consider the Riesz potential
\[
 (T_\lambda f)(t)=\int_0^1\frac{f(u)}{|t-u|^\lambda}\,du,
 \qquad t\in\Omega_\alpha.
\]
The supporting paper computes its $V$-set exactly:
\[
 V(T_\lambda)=
 \left\{\left(x,\frac{x-1+\lambda}{\alpha}\right):
 1-\lambda<x<
 \min\left\{1,\frac{1-\lambda}{1-\alpha}\right\}\right\}.
\]
This is a segment of slope $1/\alpha>1$.  The proof combines weak-type
estimates, interpolation, strict singularity of positive integral operators
off the diagonal, and a concentration sequence proving non-compactness.
Duality supplies slopes between $0$ and $1$.  Separate factorizations
through Rademacher and disjoint sequences produce the horizontal and vertical
cases.

\section*{Verification and scope}

The source PDF poses the question on page 14.  The supporting PDF identifies
that precise question on page 2 and states the exact realization theorem on
page 8.  Both papers use the same $V(T)=S(T)\setminus K(T)$ definition for
operators $L_\infty\to L_1$.  The existing run packet about arXiv:2001.09677
addresses a different endpoint question in that later paper, so it is not a
duplicate of this result.

\begin{thebibliography}{9}

\bibitem{HST17}
F. L. Hernández, E. M. Semenov, and P. Tradacete,
\emph{Interpolation and extrapolation of strictly singular operators between
$L_p$ spaces},
Adv. Math. 316 (2017), 667--690; arXiv:1706.08682.

\bibitem{HST20}
F. L. Hernández, E. M. Semenov, and P. Tradacete,
\emph{Strictly singular non-compact operators between $L_p$ spaces},
Rev. Mat. Iberoam. 39 (2023), no.~1, 181--200;
arXiv:2001.09677; doi:10.4171/RMI/1360.

\end{thebibliography}

\end{document}
